\section{Experimental Setup}
\label{sec:experimental_setup}

To rigorously evaluate the proposed \textit{PredCLOCK} (ML-CLOCK)
mechanism, we developed a custom simulation environment capable of
trace replay and machine learning inference. This section details the
hardware infrastructure, dataset characteristics, baseline
algorithms, and performance metrics used in our evaluation.

\subsection{Simulation Environment}
The experiments were conducted on a high-performance server designed
to handle large-scale trace simulations and parallel processing. The
specific specifications are chosen to ensure the entire trace and
cache structures reside in memory, minimizing I/O bottlenecks during
the simulation \cite{Yang2021A}.

\begin{itemize}
  \item \textbf{Hardware:} The testbed consists of an Intel Xeon
    Silver 4114 processor (40 logical cores at 3.00 GHz) and 192 GiB of RAM.
  \item \textbf{Software Stack:} The simulator is implemented in C++
    for efficiency, utilizing \textit{LibCacheSim} \cite{Yang2021A}
    as the foundational library for cache modeling and trace parsing.
  \item \textbf{Inference Engine:} For the ML-based eviction logic,
    we integrated \textit{ONNX Runtime} \cite{Tian2020Compiling}
    directly into the C++ environment. This allows the system to
    execute models trained in Python (via Keras/TensorFlow
    \cite{chollet2015keras, tensorflow2015-whitepaper}) with minimal
    latency overhead during the eviction decision process
    \cite{Jajal2024Interoperability}.
\end{itemize}

\subsection{Workload Datasets}
To ensure reproducibility and isolate the impact of access patterns,
we utilized synthetic traces following a Zipfian distribution.
Zipfian distributions accurately model real-world web traffic, where
a small fraction of content accounts for the majority of requests
\cite{Breslau1999Web}.

The datasets were generated procedurally using \textit{LibCacheSim}
utilities in the \texttt{oracleGeneral} binary format.
Table~\ref{tab:dataset_specs} summarizes the workload parameters. We
configured the Zipfian parameter ($\alpha$) to 1.0, representing a
standard skew found in web workloads. The stationary nature of these
traces allows us to evaluate the algorithm's steady-state performance
without the noise of phase changes.

\begin{table}[t]
  \caption{Synthetic Zipfian Dataset Specifications}
  \label{tab:dataset_specs}
  \centering
  \begin{tabular}{|l|l|}
    \hline
    \textbf{Parameter} & \textbf{Value} \\
    \hline
    Total Requests & 10,000,000 \\
    \hline
    Unique Objects & 1,000,000 \\
    \hline
    Zipf Parameter ($\alpha$) & 1.0 \\
    \hline
    Trace Format & oracleGeneral (Binary) \\
    \hline
    Data per Entry & 16 bytes (Time + ID + Size) \\
    \hline
    Distribution Type & Stationary \\
    \hline
  \end{tabular}
\end{table}

\subsection{Baselines for Comparison}
We compare the performance of our proposed \textit{PredCLOCK} against
three distinct baselines to evaluate its efficiency and effectiveness:

\begin{enumerate}
  \item \textbf{Standard CLOCK:} The unmodified CLOCK algorithm
    serves as the primary baseline. It operates purely on recency
    using a circular buffer and reference bits. It is known to suffer
    from scan resistance issues, where one-time access objects flush
    useful data.
  \item \textbf{Offline CLOCK (Theoretical Upper Bound):} This
    algorithm acts as an oracle with perfect knowledge of future
    accesses. It performs an initial scan of the trace to identify
    promotions that result in hits. If a promotion is deemed "Wasted"
    (the object is evicted before being hit), Offline CLOCK
    suppresses the promotion. This represents the theoretical maximum
    efficiency achievable \cite{Belady1966A}.
  \item \textbf{Optimal Zipf Promotion:} A heuristic-based oracle
    that promotes objects based solely on their global popularity
    ranking within the Zipfian distribution. This baseline highlights
    the value of global frequency knowledge versus local access patterns.
\end{enumerate}

\subsection{Evaluation Metrics}
To measure the effectiveness of the cache replacement strategies, we
focus on two key metrics:
\begin{itemize}
  \item \textbf{Miss Ratio:} The fraction of requests that are not
    served from the cache. A lower miss ratio indicates better cache utility.
  \item \textbf{Promotion Rate:} The total number of items promoted
    from the probationary state to the protected state (or inserted
    into the main cache). Reducing unnecessary promotions is critical
    for minimizing memory writes, which is particularly important for
    extending the lifespan of SSD-based caches \cite{Sun2022NCache,
    Wang2018Efficient}.
\end{itemize}

Experiments were conducted across various cache sizes, ranging from
1\% to 40\% of the total unique object footprint, to analyze
performance under different levels of cache pressure.
